% !TeX root = ../../portfolio.tex
% O exemplo de bloco P01 pode ser repetido para todos os projetos do percurso.
\chapter{Introdução}
Apresente o percurso, o contexto acadêmico, os objetivos do portfólio e o
período abrangido. Explique a questão que orienta a leitura do conjunto.
Considere todos os projetos pertencentes ao percurso declarado, de acordo
com o curso e o manual em elaboração \citeind{manualundf}. O documento deve
relacionar experiências, resultados e contribuição individual.

\chapter{Organização do portfólio}
Delimite o conjunto de projetos e explique como recuperou as informações:
relatórios parciais, arquivos, versões, imagens, registros e referências.
Identifique lacunas de memória ou de evidência. Use códigos de projeto e
evidência que permitam relacionar o relato aos materiais. Não fixe uma
quantidade universal de projetos nem apresente o exemplo como regra do curso.

\chapter{Visão geral dos projetos}
Registre todos os projetos do percurso informado. Acrescente uma linha por
projeto, identificando período, contexto, sua atuação e evidências. A linha
abaixo é um ponto de preenchimento, não um limite ou quantidade mínima.
\begin{undftabela}[H]{Inventário dos projetos do portfólio}{inventario-portfolio}[Elaboração do autor a partir dos registros dos projetos]
\begin{tabular}{@{}p{0.1\linewidth}p{0.25\linewidth}p{0.25\linewidth}p{0.3\linewidth}@{}}
\hline
\textbf{ID} & \textbf{Projeto e período} & \textbf{Minha contribuição} & \textbf{Evidências recuperáveis} \\
\hline
P01 & Nome e período & Atuação efetiva & Códigos e localização \\
\hline
\end{tabular}
\end{undftabela}
Explique como os projetos se distribuem no percurso. Diferencie participação
individual e produção coletiva; mencione resultados não alcançados e materiais
não recuperados quando forem relevantes à interpretação.

\chapter{Desenvolvimento dos projetos}
\section{P01 --- Nome do projeto}
Apresente contexto, problema e objetivos. Sintetize o que realizou, métodos,
ferramentas e decisões, distinguindo sua atuação da dos colegas. Discuta
entregas, resultados, limitações, dificuldades e aprendizado com referência
às evidências. Recupere o relatório parcial como fonte de memória e reescreva
o relato para a análise final. Repita a estrutura para todos os demais projetos,
adaptando o detalhamento à contribuição de cada um para o conjunto.

\chapter{Análise integradora da trajetória}
Relacione os projetos: conhecimentos mobilizados, mudanças nas decisões,
competências desenvolvidas, dificuldades recorrentes, diferenças de contexto
e resultados. Sustente essa análise com exemplos e evidências identificadas.
Explique o que aprendeu ao comparar as experiências e reconheça os limites
da documentação e da avaliação. A análise deve ultrapassar a sequência de
resumos isolados e não atribuir progresso pessoal sem justificativa.

\chapter{Conclusão}
Responda aos objetivos do portfólio e sintetize a contribuição do conjunto,
a trajetória do estudante e as limitações. Apresente recomendações coerentes
com o que foi documentado. Preserve a distinção entre resultados observados,
interpretações e possibilidades futuras. Referências e materiais complementares
seguem os elementos acadêmicos já adotados para o TCC.
